\section{第一部分实验：三档模型与三个轴}\label{sec:resultsA}

\subsection{三档模型与数据}

为了把“结构先验”与“学习”分开，我们设三档模型，共享同一个任务（给定图与源点 $s$，预测到各点的距离）：

\begin{enumerate}[leftmargin=1.8em]
\item \textbf{L0（恒等 min-sum，0 参数）}：定理~\ref{thm:equiv} 的模型，取 $g(w)=w$、不训练、推理时迭代 128 轮。它\textbf{就是} Bellman-Ford。
\item \textbf{L1（训练后的 min-sum，2 参数）}：边权变换换成 $g_\theta(w)=\mathrm{softplus}(a)\,w+b$，用监督损失训练 40 轮，推理同样迭代 128 轮。
\item \textbf{L2（mean 聚合 GNN，$104{,}321$ 参数）}：黑盒基线，每层独立参数，隐层 64，深度 8。
\end{enumerate}

必须强调：L1 与 L2 的差异是\textbf{结构先验 + 迭代预算 + 参数量}三者的联合差异，不是受控消融。干净可比的只有 L0 与 L1——它们同为 128 轮、只差 2 个参数。

训练集是 200 张 ER 图（$n\in[12,50]$，平均度 3.0），每张取 2 个源点。测试集六个：同分布 ER、跨生成器（BA / 网格 / SBM）、规模外推（$n=100,200$）。评测目标 $(s,t)$ 取 $s$ 的跳数最远点（早期版本按加权最远选，12.4\% 的实例只差 1--2 跳，等于把深度这个变量废掉了）。

\subsection{准确率：均值会骗人}

\begin{table}[H]
\centering
\caption{六个测试集上的平均绝对误差（8 个种子）}
\label{tab:acc}
\small
\setlength{\tabcolsep}{4pt}
\begin{tabular}{lcccccc}
\toprule
模型 & ER & BA & 网格 & SBM & $n{=}100$ & $n{=}200$ \\
\midrule
L0（恒等，0 参数） & 0.000 & 0.000 & 0.000 & 0.000 & 0.000 & 0.000 \\
L1（2 参数，$\times128$） & 0.024 & 0.019 & 0.033 & 0.047 & 0.037 & 0.041 \\
L2（$10^5$ 参数，8 层） & 11.77 & 8.83 & 15.35 & 28.76 & 18.95 & 19.66 \\
\bottomrule
\end{tabular}
\end{table}

\begin{table}[H]
\centering
\caption{命中率（$|\hat d-d|\le0.5$）与最坏误差（同分布 ER / 规模外推 $n{=}200$）}
\label{tab:hit}
\small
\setlength{\tabcolsep}{5pt}
\begin{tabular}{lcccc}
\toprule
\multirow{2}{*}{模型} & \multicolumn{2}{c}{ER} & \multicolumn{2}{c}{$n{=}200$} \\
\cmidrule(lr){2-3}\cmidrule(lr){4-5}
 & 命中率 & 最坏误差 & 命中率 & 最坏误差 \\
\midrule
L0 & 1.00 & 0.00 & 1.00 & 0.00 \\
L1 & 1.00 & 0.10 & 1.00 & 0.19 \\
L2 & 0.07 & 51.0 & 0.02 & 121.0 \\
\bottomrule
\end{tabular}
\end{table}

表~\ref{tab:acc}、\ref{tab:hit} 里最该看的是\textbf{命中率}与\textbf{最坏误差}，不是 MAE。黑盒的 MAE 是 11.77，听起来像“差得不多”；但它的命中率只有 0.07——\textbf{93\% 的节点误差超过 0.5}，最坏误差 121.0。均值把灾难性错误平均掉了。

还有一件必须说明的事：L2 的 11.77 是 8 个种子的均值，其中 \textbf{6/8 塌缩成了常数预测}（同一个 MAE $=14.92$），剩下两个是 2.46 与 2.19。用中位数表示是 14.92，用最好的种子表示是 2.19。所以 L2 与 L1 的比较不是“两个稳定模型之间的比较”，而是一个双峰分布与一个单峰分布的比较。

同深度对照（都只跑 8 层，不给结构臂“迭代到收敛”这个旋钮）时，结论会明显变弱：

\begin{table}[H]
\centering
\caption{同深度（8 层）对照：结构臂并非全面占优（括号内为 Mann-Whitney 双侧 $p$）}
\label{tab:same}
\small
\setlength{\tabcolsep}{4pt}
\begin{tabular}{lcccccc}
\toprule
模型 & ER & BA & 网格 & SBM & $n{=}100$ & $n{=}200$ \\
\midrule
L1 参数 @8 层 & 0.391 & 0.019 & 9.62 & 35.38 & 11.66 & 12.46 \\
L2 @8 层 & 11.77 & 8.83 & 15.35 & \textbf{28.76} & 18.95 & 19.66 \\
\midrule
$p$ 值 & $6.8\times10^{-4}$ & $6.8\times10^{-4}$ & 0.095 & $6.8\times10^{-4}$ & 0.095 & 0.095 \\
\bottomrule
\end{tabular}
\end{table}

表~\ref{tab:same} 说明：不给迭代预算时，结构臂在 ER/BA 上大幅领先，在网格与规模外推上\emph{不显著}（$p=0.095$），在 SBM 上\textbf{反而更差}（35.38 对 28.76，$p<0.001$）。原因是 SBM 的直径大于 8，8 层截断直接算不到——这正是命题~\ref{prop:depth} 的截断效应，不是“学习失败”。

\subsection{训练是负收益}

表~\ref{tab:acc} 里最反直觉的一行是 L0 与 L1：两者同为 128 轮、架构相同，只差 2 个可学习参数。未训练的 L0 误差是 $\sim10^{-6}$（float32 舍入级），训练后的 L1 却退化到 $0.024$。

先把单位说清楚。模型在归一化空间工作（边权除以 10），学到的是 $\tilde g(\tilde w)=0.9988\tilde w+0.0015$；换回真实边权即
\begin{equation}
g(w)=0.9988w+0.0146,\qquad g(w)-w=0.0146-0.0012w ,
\end{equation}
在 $w\in[1,10]$ 上每条边被高估 $+0.003\sim+0.013$。所以 $g(5)=5.0088$ 而不是 $4.9955$——那个容易被误算的值来自把归一化空间的截距当成了真实空间的截距。

这里要修正一个直觉：误差\emph{不是}“128 轮里每轮多算 0.0015、累积成 0.19”，而是\textbf{每条边被系统性高估}，迭代只是把这个每边偏差汇总成路径长度。最短路通常 3--5 条边，$3\times0.013\approx0.04$、$5\times0.003\approx0.015$，正好覆盖实测的 $0.024$。

为什么没学到 1？初始化 $a=b=0$ 对应 $g=0.693w$，训练把它从 0.693 拉到 0.9988，说明参数确实在学；但精确解在 $a=\log(e-1)\approx0.5413,\ b=0$，优化停在 $a=0.5395$，差约 0.1\%。监督信号只覆盖“8 层内可达”的节点，对 $(a,b)$ 的约束不足以把最后这点差距定出来（40 轮内损失从 0.038 降到 $10^{-4}$，曲线见仓库 \texttt{results/black.json}）。结论要限定：\textbf{在这个 2 参数化、8 层监督、40 轮、单一学习率的设置下}，训练能逼近但不能精确达到精确解；而精确解本来就在假设类里，不需要学。

\begin{figure}[H]
\centering
\includegraphics[width=0.72\textwidth]{fail_g_bias.png}
\caption{左：训练后的边权变换与恒等的偏差 $g(w)-w$，在 $w\in[1,10]$ 上为 $+0.003\sim+0.013$。右：黑盒 L2 在 60 张测试图上的预测——横轴是真值（$0$–$60$），纵轴是预测，全部挤在负值附近。}
\label{fig:gbias}
\end{figure}

\subsection{鲁棒性：不敏感不等于稳健}

\begin{table}[H]
\centering
\caption{扰动下的 MAE（权重噪声 20\% / 删边 10\%；另列任务本身的变化量）}
\label{tab:robustexp}
\begin{tabular}{lccc}
\toprule
模型 & 干净 & 权重噪声后 & 删边后 \\
\midrule
L0 & $\sim10^{-6}$ & $\sim10^{-6}$ & $\sim10^{-6}$ \\
L1 & 0.035 & 0.035 & 0.036 \\
L2 & 22.93 & 22.83 & 24.43 \\
\midrule
任务本身的变化 & --- & 5.8\% & 9.6\% \\
\bottomrule
\end{tabular}
\end{table}

表~\ref{tab:robustexp} 里有两件事值得说。第一，L0 的误差在扰动下仍然是浮点零——它本来就是在扰动后的图上重新算，这正是“白盒的鲁棒性等于可证明性”的意思。第二，L2 的误差从 22.93 变成 22.83，几乎不动。\textbf{这不能读成“稳健”}：它的误差已经接近常数预测器的水平，扰动改变不了它。真正的稳健应该是“误差小且随扰动可控”，L2 两条都不满足。

\subsection{可解释性：保真度会被常数预测器骗过}

先看一个“漂亮”的数字：L2 的 Fidelity$^+$ 是 \textbf{0.978}——只保留它的 top-$k$ 解释边，预测几乎不变。但同一个模型在解释实验的 28 张图上，预测与真值的相关系数只有 \textbf{0.13}、标准差仅为真值的 \textbf{2.9\%}；在失败案例用的 60 张图上相关系数是 $\mathbf{-0.10}$、负预测比例 100\%。一个几乎不看输入的模型，删掉任何边预测都不会怎么变——保真度自然高。\textbf{高保真度在这里度量的是模型的迟钝，不是解释的质量。}

因此跨臂比较换成三个与任务绑定的指标：

\begin{table}[H]
\centering
\caption{解释与任务的关系（28/30 张图；2 张因 $s=t$ 退化被丢弃）}
\label{tab:interp}
\small
\setlength{\tabcolsep}{4pt}
\begin{tabular}{lccccc}
\toprule
模型 & 解释边精度 & 成路率 & 放权 $\Delta$ & Fid$^+$ & 随机 Fid$^+$ \\
\midrule
L0 & 1.00 & 100\% & 3.41 & 1.000 & $-4452$ \\
L1 & 1.00 & 100\% & 3.41 & 1.000 & $-4446$ \\
L2 & 0.50 & 11\% & 4.17 & 0.978 & 0.29 \\
\bottomrule
\end{tabular}
\end{table}

其中“成路率”指解释边能连成一条 $s$--$t$ 路的样本比例，“放权 $\Delta$”指把这些边的权重放大 10 倍后\emph{真实}最短距离的相对变化。

结论要分两句话说。\textbf{其一}，黑盒的梯度显著性对\emph{模型}是忠实的：Fidelity$^+$（0.978）远高于随机选 $k$ 条边（0.29）。\textbf{其二}，它在“解释边精度”与“成路率”两项上与任务脱节：只有 11\% 的样本能连成 $s$--$t$ 路，精度 0.50（随机基线 0.00），而结构臂是 100\% 与 1.00。

还有一条观察，我原本想写成“黑盒解释无关紧要”，被数据否掉了：把解释边权重放大 10 倍，黑盒的真实距离变化是 4.17，比结构臂的 3.41 还大。原因是显著性在并列分数上按边号取 top-$k$，选中的多是与源点相连的边，恰好落在最短路的头几跳上。\textbf{“有影响”不等于“解释对了”}。

\subsection{失败案例}

只看最大误差样本容易变成挑樱桃，所以先给分布：60 张测试图上，L2 的中位误差 24.56、90 分位 33.61、\textbf{负预测比例 100\%}（预测区间 $[-2.69,-1.19]$），预测与真值的相关系数 $r=-0.10$；同一批图上，把 L1 的\emph{参数}只用 8 层推理（与 L2 同深度），中位误差 0.04、均值 0.391、不可达比例 2\%。

\begin{figure}[H]
\centering
\includegraphics[width=0.42\textwidth]{fail_depth_limit.png}
\caption{失败案例：最短路径 5 跳，而 3 层恒等 min-sum 只有 3 跳感受野，于是报“不可达”。这是\emph{构造例}（既不是 L1 也不是 L2），用于演示机制；在 60 张自然图上，8 层模型的触发率是 2\%。}
\label{fig:depthfail}
\end{figure}

三类失败对应三种机制：（1）黑盒没有值域约束，输出的距离可以为负，而距离在定义上非负，这不是某个样本的异常而是整批图的常态；（2）结构臂受深度限制（图~\ref{fig:depthfail}），$K$ 层只能看 $K$ 跳，超出即不可达，训练救不了；（3）训练的微小偏差被迭代放大（图~\ref{fig:gbias} 左）。

需要说明口径：本节用 60 张图（种子 3000--3059），第~\ref{sec:resultsA}~节的 $11.77$ 用 40 张（3000--3039），两个“最坏误差”（51.0 与 58.9）不可直接并列。
